\documentclass[UTF8]{article}
\usepackage{amsmath, amssymb, amsthm}
\usepackage{geometry}
\geometry{a4paper, margin=1in}
\usepackage{ctex} % 支持中文
\usepackage{xcolor}

\title{一元二次不等式解答题}
\author{}
\date{}

\begin{document}
	
	
	\section*{例 1}
	\textbf{题目：}求不等式 \(x^2 - x + m(1-m) > 0\) 的解集。
	
	\vspace{0.5cm}
	\textbf{解答：}
	
	将不等式左边因式分解：
	\[
	x^2 - x + m(1-m) = (x-m)(x + m - 1)
	\]
	原不等式等价于：
	\[
	(x-m)[x - (1-m)] > 0
	\]
	对应方程的两个根为 \(x_1 = m\) 和 \(x_2 = 1 - m\)。
	
	比较两个根的大小：
	\begin{itemize}
		\item 令 \(m = 1 - m\)，解得 \(m = \frac{1}{2}\)。
		\item 令 \(m > 1 - m\)，解得 \(m > \frac{1}{2}\)。
		\item 令 \(m < 1 - m\)，解得 \(m < \frac{1}{2}\)。
	\end{itemize}
	
	根据根的大小关系，分情况写出解集（大于取两边）：
	
	\begin{enumerate}
		\item \textbf{当 \(m > \frac{1}{2}\) 时，\(m > 1-m\)：}
		解集为 \(\{x \mid x < 1-m \text{ 或 } x > m\}\)。
		
		\item \textbf{当 \(m = \frac{1}{2}\) 时，\(m = 1-m = \frac{1}{2}\)：}
		不等式变为 \((x - \frac{1}{2})^2 > 0\)。
		解集为 \(\{x \mid x \neq \frac{1}{2}\}\)。
		
		\item \textbf{当 \(m < \frac{1}{2}\) 时，\(m < 1-m\)：}
		解集为 \(\{x \mid x < m \text{ 或 } x > 1-m\}\)。
	\end{enumerate}
	
	\newpage
	
	\section*{例 2}
	\textbf{题目：}求不等式 \(x^2 - 2mx + (m+2) < 0\) 的解集。
	
	\vspace{0.5cm}
	\textbf{解答：}
	
	考虑对应的一元二次方程 \(x^2 - 2mx + (m+2) = 0\)。
	其判别式 \(\Delta = (-2m)^2 - 4 \times 1 \times (m+2) = 4m^2 - 4m - 8 = 4(m^2 - m - 2) = 4(m-2)(m+1)\)。
	
	根据判别式的符号，我们分情况讨论：
	
	\begin{enumerate}
		\item \textbf{当 \(\Delta < 0\) 时，即 \(-1 < m < 2\) 时：}
		抛物线开口向上且与 \(x\) 轴无交点，不等式 \(x^2 - 2mx + (m+2) < 0\) 无解。
		此时解集为 \(\emptyset\)。
		
		\item \textbf{当 \(\Delta = 0\) 时，即 \(m = -1\) 或 \(m = 2\) 时：}
		方程有两个相等的实根。
		\begin{itemize}
			\item 若 \(m = -1\)，方程为 \(x^2 + 2x + 1 = 0\)，即 \((x+1)^2 = 0\)，根为 \(x = -1\)。不等式 \((x+1)^2 < 0\) 无解，解集为 \(\emptyset\)。
			\item 若 \(m = 2\)，方程为 \(x^2 - 4x + 4 = 0\)，即 \((x-2)^2 = 0\)，根为 \(x = 2\)。不等式 \((x-2)^2 < 0\) 无解，解集为 \(\emptyset\)。
		\end{itemize}
		
		\item \textbf{当 \(\Delta > 0\) 时，即 \(m < -1\) 或 \(m > 2\) 时：}
		方程有两个不相等的实根，由求根公式得：
		\[
		x = \frac{2m \pm \sqrt{4(m-2)(m+1)}}{2} = m \pm \sqrt{m^2 - m - 2}
		\]
		因为抛物线开口向上，不等式 \(< 0\) 的解集在两根之间。
		此时解集为 \(\{x \mid m - \sqrt{m^2 - m - 2} < x < m + \sqrt{m^2 - m - 2}\}\)。
	\end{enumerate}
	
	\textbf{综上所述：}
	\begin{itemize}
		\item 当 \(-1 \le m \le 2\) 时，解集为 \(\emptyset\)；
		\item 当 \(m < -1\) 或 \(m > 2\) 时，解集为 \((m - \sqrt{m^2 - m - 2}, m + \sqrt{m^2 - m - 2})\)。
	\end{itemize}
	
	\newpage
	

	
	\section*{例 3}
	\textbf{题目：}已知不等式 \(x^2 + mx + n < 0\) 的解集为 \(\{x \mid 1 < x < 4\}\)，
	(1) 求 \(m\) 和 \(n\) 的值。
	(2) 若 \(x^2 + mx + n \le ax\) 对一切 \(x > 0\) 恒成立，求实数 \(a\) 的取值范围。
	
	\vspace{0.5cm}
	\textbf{解答：}
	
	\textbf{(1) 求 \(m\) 和 \(n\) 的值：}
	
	由题意知，不等式 \(x^2 + mx + n < 0\) 的解集为 \((1, 4)\)。
	这说明二次方程 \(x^2 + mx + n = 0\) 的两个根分别为 \(x_1 = 1\) 和 \(x_2 = 4\)。
	
	根据韦达定理（根与系数的关系）：
	\begin{align*}
		x_1 + x_2 &= -m \implies 1 + 4 = -m \implies m = -5 \\
		x_1 \cdot x_2 &= n \implies 1 \times 4 = n \implies n = 4
	\end{align*}
	所以，\(m = -5\)，\(n = 4\)。
	
	\vspace{0.5cm}
	\textbf{(2) 求实数 \(a\) 的取值范围：}
	
	由 (1) 知，不等式为 \(x^2 - 5x + 4 \le ax\)。
	整理得：\(x^2 - (5+a)x + 4 \le 0\)。
	该不等式对一切 \(x > 0\) 恒成立。
	
	由于 \(x > 0\)，我们可以将不等式两边同时除以 \(x\)，得：
	\[
	x - (5+a) + \frac{4}{x} \le 0
	\]
	\[
	x + \frac{4}{x} \le 5 + a
	\]
	要使该不等式对一切 \(x > 0\) 恒成立，只需 \(5+a\) 大于等于函数 \(f(x) = x + \frac{4}{x}\) 在 \(x > 0\) 时的最小值。
	
	根据基本不等式（AM-GM）：
	\[
	x + \frac{4}{x} \ge 2\sqrt{x \cdot \frac{4}{x}} = 2\sqrt{4} = 4
	\]
	当且仅当 \(x = \frac{4}{x}\)，即 \(x = 2\) 时（满足 \(x>0\)），等号成立。
	所以，\(f(x)\) 的最小值为 \(4\)。
	
	因此，我们只需要：
	\[
	5 + a \ge 4
	\]
	\[
	a \ge -1
	\]
	
	所以，实数 \(a\) 的取值范围是 \([-1, +\infty)\)。
	
\end{document}